\booksection{נספח פרק 11: ידע טכני מקצועי בצילום ווידאו}

\chapterintro{כדי באמת להבין את סודות הצילום והווידאו בטלפון, אתה צריך להכיר ידע טכני מקצועי על צילום. ידע זה לא רק יעזור לך להשתמש טוב יותר במצלמת הטלפון, אלא גם יעניק לך שיקול דעת טוב יותר בבחירת מכשירים.}

\booksubsection{רזולוציה, קצב פריימים וקצב נתונים: שלושת הגורמים המרכזיים באיכות וידאו}

\subsubsection*{720P (HD)}

: 1280×720 פיקסלים, רמת בהירות בסיסית.

\begin{itemize}
  \item \textbf{1080P (Full HD)}: 1920×1080 פיקסלים, הרזולוציה הנפוצה ביותר כיום.
  \item \textbf{2K (QHD)}: 2560×1440 פיקסלים, חד יותר מ-1080P.
  \item \textbf{4K (UHD)}: 3840×2160 פיקסלים, אולטרה-HD, עשיר בפרטים.
  \item \textbf{8K}: 7680×4320 פיקסלים, תקן העתיד, תמיכת מכשירים מוגבלת כיום.
  \item \textbf{ככל שהרזולוציה גבוהה יותר, התמונה חדה יותר, אבל הקובץ גדול יותר ודורש יותר משאבי אחסון ועיבוד.}קצב פריימים (Frame Rate):
\end{itemize}

קצב הפריימים קובע את חלקות התנועה, מבוטא בפריימים לשנייה (fps).

\subsubsection*{24fps}

: קצב פריימים סטנדרטי בקולנוע, בעל “תחושת קולנוע” ייחודית.

\begin{itemize}
  \item \textbf{30fps}: נפוץ בטלוויזיה ובווידאו אינטרנטי, חלק וטבעי.
  \item \textbf{60fps}: קצב פריימים גבוה, מתאים לסצנות ספורט והקלטת משחקים, חלק מאוד.
  \item \textbf{120fps/240fps}: קצב פריימים אולטרה-גבוה, ניתן לשימוש לווידאו בהילוך איטי.
  \item \textbf{ככל שקצב הפריימים גבוה יותר, התנועה חלקה יותר ואפקט ההילוך האיטי טוב יותר, אבל הקובץ גדול יותר.}קצב נתונים (Bitrate):
\end{itemize}

קצב הנתונים קובע את צפיפות המידע בווידאו, מבוטא בכמות נתונים לשנייה (Mbps).

\subsubsection*{ככל שקצב הנתונים גבוה יותר, איכות התמונה טובה יותר ונשמרים יותר פרטים, אבל הקובץ גדול יותר.}

קצב נתונים נמוך מדי גורם להופעת סימני דחיסה כמו פסיפס וכתמי צבע בתמונה.

\begin{itemize}
  \item צילום וידאו מקצועי דורש בדרך כלל קצב נתונים גבוה כדי לשמור מרחב עיבוד גדול יותר.
  \item 2. פורמטי קידוד וידאו: H.264, H.265 ו-ProRes
  \item H.264 (AVC):
\end{itemize}

\booksubsection{פורמטי קידוד וידאו: H.264, H.265 ו-ProRes}

\subsubsection*{: פורמט הקידוד הנפוץ ביותר, תאימות מעולה.}

\begin{itemize}
  \item \textbf{דרישות אביזרים}: כמעט כל המכשירים יכולים לנגן, יעילות דחיסה סבירה.
  \item \textbf{חסרונות}: בהשוואה לפורמטים חדשים, יעילות דחיסה נמוכה יותר, קבצים גדולים יותר.
  \item \textbf{תרחישי שימוש}: שיתוף וידאו באינטרנט, צילום יומיומי למשתמש רגיל.
  \item \textbf{3. HDR: טכנולוגיית טווח דינמי גבוה}מאפיינים
\end{itemize}

\subsubsection*{: גרסה משודרגת של H.264, יעילות דחיסה משופרת בכ-50\%.}

\begin{itemize}
  \item \textbf{דרישות אביזרים}: באותה איכות תמונה הקובץ קטן יותר, תומך ברזולוציה וקצב פריימים גבוהים יותר.
  \item \textbf{חסרונות}: קידוד ופענוח דורשים מעבד חזק יותר, חלק מהמכשירים הישנים לא תואמים.
  \item \textbf{תרחישי שימוש}: צילום וידאו 4K, תרחישים עם שטח אחסון מוגבל.
  \item \textbf{3. HDR: טכנולוגיית טווח דינמי גבוה}מאפיינים
\end{itemize}

\subsubsection*{: קידוד וידאו מקצועי שפותח על ידי Apple, דחיסה כמעט ללא אובדן.}

\begin{itemize}
  \item \textbf{דרישות אביזרים}: איכות תמונה גבוהה במיוחד, מרחב עיבוד גדול, צבע מדויק.
  \item \textbf{חסרונות}: קבצים גדולים מאוד, דורש נפח אחסון רב, עיבוד דורש ציוד מקצועי.
  \item \textbf{תרחישי שימוש}: הפקת וידאו מקצועית, צילום ברמת קולנוע.
  \item \textbf{3. HDR: טכנולוגיית טווח דינמי גבוה}מאפיינים
\end{itemize}

\subsubsection*{: משלב את היתרונות של פורמט ProRes ופורמט RAW.}

\begin{itemize}
  \item \textbf{דרישות אביזרים}: נתוני חיישן גולמיים לא מעובדים, שומרים על כמות המידע המרבית.
  \item \textbf{יתרונות}: מרחב התאמה עצום לאחר הצילום, ניתן לשנות פרמטרים כמו חשיפה ואיזון לבן ללא אובדן איכות.
  \item \textbf{חסרונות}: קבצים גדולים במיוחד, דורש תוכנת עיבוד מקצועית.
  \item \textbf{תרחישי שימוש}: הפקה מקצועית יוקרתית. סדרת iPhone 15 Pro תומכת בצילום וידאו Apple ProRes (עד 4K 60fps), אך פורמט ProRes RAW דורש כרגע ציוד הקלטה חיצוני, צילום ישיר בטלפון עדיין לא נתמך.
  \item \textbf{3. HDR: טכנולוגיית טווח דינמי גבוה}:הפקה מקצועית. iPhone 15 Pro ואילך תומך בצילום וידאו Apple ProRes (עד 4K 60fps), iPhone 17 Pro תומך ב-ProRes RAW בתוך המכשיר.
\end{itemize}

\booksubsection{HDR: טכנולוגיית טווח דינמי גבוה}

: תקן HDR הבסיסי ביותר, תאימות טובה.

\begin{itemize}
  \item \textbf{HDR10+}: תקן שמקודם על ידי Samsung, תומך במטא-דאטה דינמי.
  \item \textbf{Dolby Vision}: תקן HDR עילי שפותח על ידי Dolby, איכות תמונה מיטבית, אך דורש רישוי.
  \item \textbf{HLG}: תקן HDR ברמת שידור, תואם למסכים רגילים.
  \item \textbf{בעת צילום וידאו HDR, נדרש מסך תומך HDR כדי לראות את האפקט המלא.}4. מרחב צבע ועומק צבע
\end{itemize}

מרחב צבע:

\booksubsection{מרחב צבע ועומק צבע}

\subsubsection*{: מרחב צבע סטנדרטי, התאימות הטובה ביותר.}

\begin{itemize}
  \item \textbf{Display P3}: תקן מרחב צבע רחב שמקודם על ידי Apple, צבעים עשירים יותר.
  \item \textbf{Rec.2020}: מרחב הצבע הסטנדרטי לווידאו HDR, מרחב הצבע הרחב ביותר.
  \item \textbf{עומק צבע:}8-bit
\end{itemize}

\subsubsection*{: 16.7 מיליון צבעים, תקן וידאו רגיל.}

\begin{itemize}
  \item \textbf{10-bit}: 1.07 מיליארד צבעים, תקן וידאו מקצועי, מעברי צבע חלקים יותר.
  \item \textbf{12-bit}: 68 מיליארד צבעים, תקן מקצועי יוקרתי.
  \item \textbf{המלצות פרקטיות}צילום יומיומי: 1080P/4K + 30fps + קידוד H.265, איזון בין איכות תמונה לגודל קובץ.
\end{itemize}

\begin{tcbinfo}{צילום ספורט: 1080P/4K + 60fps, להבטחת חלקות.}
\begin{itemize}
  \item יצירה מקצועית: 4K + ProRes (אם נתמך), שמירת מרחב עיבוד מרבי.
  \item המלצת אחסון: וידאו בקצב נתונים גבוה תופס הרבה מקום, מומלץ לגבות באופן קבוע למחשב או לענן.
  \item צא לחקור: הגדרות הווידאו בטלפון שלך
  \item פתח את אפליקציית המצלמה ועבור למצב וידאו.
\end{itemize}

\end{tcbinfo}

\begin{tcbpractice}{בדוק את אפשרויות הרזולוציה וקצב הפריימים, נסה שילובים שונים.}
\begin{enumerate}
  \item אם נתמך, בדוק אם יש אפשרויות מתקדמות כמו HDR, ProRes.
  \item צלם מספר קטעי וידאו בהגדרות שונות והשווה איכות תמונה וגודל קובץ.
  \item פרק 12: בידור אודיו-ויזואלי וקריאה
  \item הטלפון כיום בעצם כבר החליף את נגן ה-MP3, שלט הטלוויזיה, ה-Kindle וקונסולת המשחקים הניידת. אבל למען האמת, רוב האנשים משתמשים רק בפונקציות הבסיסיות ביותר של האפליקציות האלה, ויש עוד המון דברים שימושיים שלא התגלו. בפרק הזה נדבר על איך להפוך את הבידור בטלפון לנוח ונעים יותר.
\end{enumerate}

\end{tcbpractice}

\booksubsection{ה. השוואת פורמט RAW מול HEIF/JPEG}

\subsubsection*{מהו פורמט RAW:}

פורמט RAW הוא ה“נתונים הגולמיים” של חיישן המצלמה - מידע לא מעובד ולא דחוס. כמו חומר גלם, דורש “בישול” בעיבוד מאוחר. טלפונים מצלמים כברירת מחדל JPEG או HEIF “מבושלים”, RAW דורש הפעלה מיוחדת.

\subsubsection*{השוואת RAW, HEIF ו-JPEG:}

\begin{itemize}
  \item \textbf{JPEG}:הפורמט הישן ביותר, תאימות מיטבית (כל מכשיר פותח). עומק צבע 8-bit, דחיסה מאבדת נתונים, קובץ קטן (2-5MB לתמונה). מרחב עיבוד מוגבל, הבהרת אזורים חשוכים עלולה ליצור גושי צבע.
  \item \textbf{HEIF/HEIC}:פורמט חדש שאפל מקדמת, באותה איכות הקובץ בגודל חצי מ-JPEG (1-3MB). תומך ב-10-bit ובערוץ שקיפות. חיסרון: מכשירים ישנים ו-Windows דורשים תוסף.
  \item \textbf{RAW}:נתונים גולמיים לא דחוסים, הקובץ הגדול ביותר (20-100MB לתמונה). עומק צבע 12-16bit, מרחב עיבוד מקסימלי. חיסרון: דורש תוכנה מקצועית (Lightroom, Snapseed) ולא ניתן לשתף ישירות.
\end{itemize}

\subsubsection*{למה צלמים מקצועיים מצלמים RAW:}

\begin{itemize}
  \item \textbf{איזון לבן ניתן לתיקון ללא אובדן בעיבוד מאוחר}:ב-JPEG אם הבלנס הלבן שגוי (כמו גוון צהבהב תחת נורת להט), תיקון מאבד איכות; ב-RAW ניתן להגדיר מחדש את טמפרטורת הצבע ללא כל איבוד.
  \item \textbf{סובלנות חשיפה רחבה}:ב-JPEG חשיפת יתר של 2 סטופים משאירה אזורים חשוכים, הבהרה יוצרת רעש; RAW יכול "להציל" פרטים מחשיפת יתר של 3-4 סטופים. גם לחשיפת יתר.
  \item \textbf{פרטים עשירים יותר בצללים ובבהירים}:טווח הדינמי של RAW גדול בהרבה מ-JPEG, עם גמישות רבה יותר בבהירים ובכהים.
  \item \textbf{חידוד, הפחתת רעש ותיקון עדשה – כולם נשלטים בעיבוד מאוחר}:העיבוד במצלמה הוא "חד-פעמי", RAW ניתן לעבד שוב ושוב, תמיד מהנתונים המקוריים.
\end{itemize}

\subsubsection*{איך מפעילים RAW בטלפון:}

\begin{itemize}
  \item \textbf{iPhone ProRAW}:רק iPhone 12 Pro ומעלה. הפעל "Apple ProRAW" ב[הגדרות] > [מצלמה] > [פורמט]. במצלמה יופיע כפתור "RAW", הדלק. קובץ כ-25MB לתמונה, הפעל רק בעת צורך.
  \item \textbf{Android Motion Raw / Ultra RAW}:שמות שונים לכל מותג. Xiaomi/OPPO/vivo ב"מצב מקצועי" או "עוד" במצלמה; Huawei ב"מצב מקצועי"; Samsung Galaxy S ב"מצב מקצועי". גודל קובץ 20-50MB.
  \item \textbf{המלצה כללית}:לצילום יומיומי אין צורך ב-RAW, הקובץ גדול מדי; הפעל רק בסיטואציות חשובות (נוף, פורטרט, עיבוד מאוחר).
\end{itemize}

\booksubsection{ו. עקרונות צילום חישובי}

צילום בטלפון מודרני כבר לא יכול בלי “צילום חישובי” - יצירת תמונות על ידי אלגוריתמים ולא רק עדשות וחיישנים. הבנת העקרונות מסבירה מה הטלפון עושה ברגע לחיצת הצילום.

\subsubsection*{סינתזה רב-פריימית:}

בלחיצת צילום אחת הטלפון מצלם מספר תמונות. האלגוריתם מיישר וממזג את הפריימים לתמונה אחת איכותית.

\begin{itemize}
  \item \textbf{מצב לילה}:צילום מרובה של תמונות חשוכות ב-1-3 שניות, יישור והדבקה - רעש מבטל זה את זה, בהירות מצטברת. תוצאה: צילום נקי באור חלש.
  \item \textbf{HDR (טווח דינמי רחב)}:צילום נפרד של תמונות חשוכות, נורמליות ובהירות, מיזוג לתמונה אחת עם פרטים בבהירים ובכהים. יעיל במיוחד נגד האור.
  \item \textbf{Deep Fusion (אפל) / Photonic Engine}:לסצנות חלשות באור, ניתוח פריימים פיקסל-פיקסל ומיזוג, שיפור מרקם ופרטים, סריגים ושיער יהיו ברורים יותר.
\end{itemize}

\subsubsection*{פילוח סמנטי:}

AI לא מעבד את כל התמונה אחת, אלא “מבין” מה בתמונה - שמיים, פנים, דשא, בניינים, בגדים - ומחיל אסטרטגיות שונות על כל אזור.

\begin{itemize}
  \item \textbf{שמיים}: הגברת כחול בנפרד, הנמכת בהירות, מניעת חשיפת יתר.
  \item \textbf{פנים}:זיהוי אזורי גוון עור, הבהרה והחלקה, בלי לגעת בפנים ובשיער.
  \item \textbf{דשא / עלים}: הגברת רוויה ירוקה, שיפור חידוד פרטים.
  \item \textbf{טשטוש רקע}:במצב פורטרט, AI מזהה את קו המתאר, מטשטש רקע ומשאיר נושא חד.
\end{itemize}

\subsubsection*{יישומי למידה עמוקה:}

\begin{itemize}
  \item \textbf{זיהוי סצנה ב-AI}:המצלמה מזהה אוטומטית סצנות כמו "אוכל" "חיות מחמד" "לילה" "מסמך" ומתאימה פרמטרים. למשל מצב אוכל מחזק גוונים חמים ורוויה.
  \item \textbf{טשטוש דיוקן ב-AI}:עדשת טלפון לא יכולה להשיג עומק שדה רדוד כמו DSLR, במקום זאת מודל למידת מכונה מנבא את קצוות הנושא ומטשטש את הרקע. אלגוריתמים חדשים מטפלים נכון בשיער ובקצוות זכוכית.
  \item \textbf{הסרה ב-AI / הרחבת תמונה ב-AI}:בחירת עוברי אורח או חפצים בתמונה, AI "משלים" את החלק החסר. זה יישום של AI גנרטיבי בצילום.
\end{itemize}

\subsubsection*{צילום חישובי מול צילום אופטי:}

\begin{itemize}
  \item \textbf{יתרונות צילום חישובי}:נפח קטן עם איכות טובה, צילום באור חלש ונגד האור, משתמש רגיל לא צריך ידע בצילום כדי לצלם "יפה".
  \item \textbf{חסרונות צילום חישובי}:לפעמים "עיבוד יתר" - חדות יתר, עיוות צבע, שמיים לא טבעיים; מיזוג פריימים יוצר הילות כשהנושא זז; פרטים "מחושבים" ולא מתועדים.
  \item \textbf{יתרונות צילום אופטי}:חיישן גדול + עדשה בצמצם רחב לוכדים אור אמיתי, איכות טבעית, מרחב עיבוד גדול; ללא פגמים מ"השלמת" אלגוריתם.
  \item \textbf{מגמות}:טלפוני דגל מודרניים הם שילוב "אופטי+חישובי" - חיישן גדול למידע אמיתי, אלגוריתם לאופטימיזציה. הם לא מנוגדים.
\end{itemize}

\booksubsection{ז. סוגי עדשות ומרחקי מוקד}

טלפוני דגל מודרניים מצוידים ב-3 עד 4 עדשות, כל אחת לטווח מוקד שונה. הבנת הטווח עוזרת לבחור את העדשה הנכונה.

\subsubsection*{עדשות נפוצות ומרחק מוקד שקול:}

\begin{itemize}
  \item \textbf{אולטרה רחבה (13-24mm שקול)}:שדה הראייה הרחב ביותר, מתאים לנוף, בניינים, חללי פנים. מצלם מרחבים גדולים אך עיוות בקצוות (קווים ישרים מתעקמים), פנים נמתחות בפורטרט. מתאים ל"מרחב גדול", לא ל"תקריב".
  \item \textbf{עדשה ראשית (24-28mm שקול)}:העדשה הנפוצה ביותר, האיכות והצמצם הטובים ביותר. מאזן שדה ראייה ואיכות. "1x" ברוב הטלפונים היא העדשה הראשית.
  \item \textbf{טלה (70-200mm שקול)}:לצילום רחוק ופורטרט. טלפוטו מדחס תחושת מרחב (הרים רחוקים ואנשים קרובים נראים קרובים יותר), פורטרט עם פרופורציות טבעיות וטשטוש רקע.
  \item \textbf{טלה פריסקופי (200mm+ שקול)}:תכנון אופטי מבוסס פריזמה, משיג טווח ארוך יותר בעובי הטלפון. מצלם פרטים מרחוק, מתאים להופעות, ספורט וחיות בר. רק בדגמי דגל.
\end{itemize}

\subsubsection*{פשרה בין גודל חיישן לגודל פיקסל:}

החיישן הוא ה“בסיס”, קובע כמה אור נקלט. בתעשייה אומרים “בסיס גדול יותר מנצח” - באותו מספר פיקסלים, חיישן גדול יותר = פיקסלים גדולים יותר = יותר אור = איכות טובה יותר.

\begin{itemize}
  \item \textbf{חיישן גדול}:קולט הרבה אור, איכות לילה טובה, טווח דינמי רחב. חיסרון: הטלפון עבה יותר (בליטת עדשה), יקר.
  \item \textbf{פיקסלים רבים, חיישן קטן}:באור יום מספיק, פיקסלים גבוהים מצלמים חד יותר, נוח לחיתוך. אך באור חלש פחות אור לפיקסל, יותר רעש.
  \item \textbf{בחירות דגל מובילות}:העדשה הראשית בדרך כלל 1/1.3 עד 1 אינץ' (חלק כבר 1 אינץ'), כ-50 מיליון פיקסלים - נקודת האיזון בין איכות ונפח.
\end{itemize}

\subsubsection*{טכנולוגיית מיזוג פיקסלים (Pixel Binning):}

כדי לשלב פיקסלים גבוהים ביום ופיקסלים גדולים בלילה, יצרנים משתמשים ב“מיזוג פיקסלים” - מאחדים מספר פיקסלים קטנים לפיקסל גדול אחד.

\begin{itemize}
  \item \textbf{פיקסלים מרובעים (Quad Bayer)}:חיישן 50 מיליון פיקסל, פלט 12.5 מיליון. כל פיקסל פלט מורכב מ-4 פיקסלים, אור פי 4, איכות נקייה יותר.
  \item \textbf{פיקסלים תשעוניים (Nonacell)}:לחיישנים מעל 100 מיליון פיקסל, פלט כ-1100-1200 מיליון. אותו עיקרון, יותר אור.
  \item \textbf{אפקט מעשי}:במצב צילום רגיל, מקבלים תמונה ממוזגת בפיקסלים נמוכים - איכות טובה, קובץ קטן. לפרטים מקסימליים (מסמך, רחוק), ניתן לעבור ל"מצב פיקסלים גבוה".
\end{itemize}

\booksubsection{ח. מצבי צילום מיוחדים}

מלבד צילום יומיומי, יש מצבים מיוחדים שיכולים לצלם מה שאי אפשר בדרך כלל.

\subsubsection*{מצב פנורמה:}

\begin{itemize}
  \item בצילום, הזז את הטלפון לאט וביציבות לרוחב לפי החץ, לא מהר - מהרות גורמת לשגיאות חיבור.
  \item שמור על יציבות אנכית, רק הזז אופקית. אם יש תזוזה אנכית, קו האופק יהיה גלי.
  \item היזהר עם אנשים או רכבים נעים - בחיבור יופיעו "חצי גוף" או "רוחות רפאים".
  \item מצב פנורמה יכול גם לצלם לאורך - צילום אנכי חיבור אופקי, מתאים לבניינים גבוהים או מפלים.
\end{itemize}

\subsubsection*{הילוך איטי וצילום טווח זמן:}

\begin{itemize}
  \item \textbf{הילוך איטי}:צילום ב-120fps או 240fps, השמעה לאט. מתאים לטיפות מים, רגעי ספורט, קפיצות חיות. דורש אור מספיק, באור חלש איכות ירודה.
  \item \textbf{טיים לאפס}:צילום מרווח של מספר תמונות לוידאו מואץ. מתאים לתנועת מכוניות, עננים, שקיעה, פריחה. נמשך דקות עד שעות, חובה להשתמש בחצובה.
  \item \textbf{טריק קטן}:בצילום עננים ב-timelapse, הנח את הטלפון, צלם 15-30 דקות, קבל כ-10 שניות של וידאו מרהיב.
\end{itemize}

\subsubsection*{מצב כוכבים / אסטרונומיה:}

\begin{itemize}
  \item יש להתרחק מזיהום אור עירוני, למצוא אזור ללא פנסי רחוב. בעיר לא ניתן לצלם את שביל החלב.
  \item חובה חצובה, יד בלתי אפשרית. הטלפון חייב להיות יציב למספר שניות.
  \item מזג אוויר צריך להיות בהיר ללא עננים, הירח לא גדול (ירח מלא מסתיר כוכבים). לילות חשוכים בתחילת או סוף החודש המומלצים.
  \item "מצב לילה" של iPhone מזהה אוטומטית שמיים כוכבים ומאריך חשיפה (עד 30 שניות); "מצב כוכבים" או "מצב אסטרונומי" ב-Android דומה. המתן בסבלנות, אל תזוז את הטלפון.
\end{itemize}

\subsubsection*{מצב מאקרו:}

\begin{itemize}
  \item לצילום פרחים, חרקים, תכשיטים, טקסט. יכול להתמקד מ-1-2 ס"מ.
  \item חלק מהטלפונים עם עדשת מאקרו ייעודית (בדרך כלל 2 או 5 מיליון פיקסל), איכות בינונית; דגמי דגל משתמשים באולטרה רחב למאקרו, איכות טובה יותר.
  \item בצילום תקריב העדשה קרובה וחוסמת אור, השתמש במנורה או פנס לתאורה.
  \item מרחק מיקוד קצר, עומק שדה רדוד, תזוזה של מילימטרים מזיזה את המוקד. יד יציבה, עדיף שהנושא לא יזוז.
\end{itemize}

\begin{tcbinfo}{המלצות מעשיות לצילום סטטי}
\begin{itemize}
  \item צילום יומיומי: HEIF/JPEG כברירת מחדל, קובץ קטן, איכות מספקת.
  \item סצנות חשובות (נוף, פורטרט, עיבוד): הפעל RAW למרחב עיבוד מקסימלי.
  \item בחר עדשה נכונה: נוף - אולטרה רחב, פורטרט - טלפוטו, יומיומי - ראשית.
  \item לילה וכוכבים: חובה לקבע את הטלפון, הפעל מצב לילה או אסטרונומי.
  \item טלפון עם חיישן גדול (1 אינץ') בעל פוטנציאל איכות גבוה, אך הרגלי צילום הם המפתח.
\end{itemize}

\end{tcbinfo}

\begin{tcbpractice}{תרגול מעשי: הכר את יכולות המצלמה שלך}
\begin{enumerate}
  \item בדוק אם הטלפון תומך ב-RAW (הגדרות > מצלמה > פורמט), הפעל וצלם, השווה גודל JPEG מול RAW.
  \item עבור בין אולטרה רחב, ראשית וטלפוטו, צלם את אותו נוף, הרגש את ההבדל בזווית ובאיכות.
  \item מצא סצנה עם מצב לילה או כוכבים, קבע את הטלפון וצלם תמונת לילה.
  \item נסה פנורמה אופקית של מרחב גדול, שמור על תנועה יציבה.
\end{enumerate}

\end{tcbpractice}
